\section{Results}
\label{sec:results}

Each subsection states a question, its design and slots for its results.
All runs use one tagged impg build, and we store its sha256 digest, version and commit next to each result \authornote{tag name, Zenodo version DOI and owner of the frozen build, decision D6 and code fix C1}.
% src: supporting/plan.md section 7 (rules for every run) and section 9 C1
Every benchmark names its index with \texttt{-i}, which forces a single index.
% src: impg src/main.rs:11277-11291 (IndexMode::Auto returns single when an index path is given)
Without \texttt{-i}, impg switches to one index per file at 100 or more files, and no test covers that mode.
% src: impg src/main.rs:11281-11287; supporting/plan.md section 9 C6 (src/multi_impg.rs:757-952, 0 tests)
A fixed-seed sampler draws every region set, and held-out sets never choose parameters (Supplementary Note S5 and Table~\ref{tab:commands}).
% src: supporting/plan.md section 7 (rules for every run)

\subsection{Tracepoint projection and CIGAR projection}
\label{sec:res-tpa}

HPRC release 2 distributes the PGGB alignments as 466 TPA files, one per haplotype \citep{HPRC2026resources,Guarracino2026tpa}.
% src: HPRC2026resources, README lines 151 and 183 (phase1/web/hpp_readme.md:151, 183)
For tracepoint input, impg rebuilds the CIGAR of the segments that overlap a query from the sequences with WFA2 \citep{Kaushan2026tracepoints,MarcoSola2021wfa}.
% src: impg src/impg.rs:1206-1362
No test compares this path with PAF input, and impg drops a hit it cannot rebuild with only a debug message.
% src: supporting/plan.md section 9 C2; impg src/impg.rs:1230-1236, 1247-1254, 1330-1338, 1436-1442, 1453-1460
We ask whether exact tracepoint projection returns the same rows as CIGAR projection of the same alignments.
We use the GRCh38 target file, 8,916,633 records, as TPA (560,062,985 bytes) and as gzipped PAF (6,532,446,867 bytes).
% src: S3 listings of release2/impg/tpas and release2/impg/pafs/all-vs-1 (phase2/s3/tpas.tsv, phase2/s3/pafs.tsv); supporting/plan.md sections 4 and 5.4 (record counts equal); data/partners.tsv line 1
We project 1,000 random 10 kb GRCh38 regions, plus 200 held out, at depth 1, 2 and unbounded, without merging.
% src: supporting/plan.md section 7 E1
We count identical sorted BEDPE rows, rows found in one output only, the largest boundary shift and dropped tracepoint hits.
% src: supporting/plan.md section 7 E1
We also run \texttt{-{}-approximate}, which serves only BED and BEDPE output, without tuning it.
% src: impg query --help (--approximate)
A second test indexes two assembly files that list each other, with and without \texttt{-{}-unidirectional}, because both files store their alignments.
% src: data/pairs_out.txt (2,966 ordered pairs stored in both files); impg src/impg.rs:1686-1712 (each non-self record is also inserted reversed); supporting/plan.md section 7 E1b
Exact tracepoint and PAF projection gave \pending{E1}{identical sorted rows, \%, per depth and set}.
The tracepoint path dropped \pending{E1}{tracepoint hits, count}.
Approximate mode shifted boundaries by \pending{E1}{median and maximum shift, bp} at \pending{E1}{median s per query, both modes}.
The reciprocal pair gave \pending{E1}{index entries per record, merged BED identity}.
Every tracepoint result below depends on this test (Supplementary Fig.~S2).

\subsection{Indexing and querying the whole release}
\label{sec:res-release}

The Shredtools authors wrote that they could not compute a comparable impg index over whole HPRCr1 assemblies or the HPRCr2 collection of N = 476 \citep{Shivakumar2026shredtools}.
% src: Shivakumar2026shredtools, references/biorxiv-2026-07-03-736354/paper.txt:285-287
The same preprint states that impg in practice runs on a sparse sample of pairs, and it benchmarked 6\% and 17\% pair sets.
% src: Shivakumar2026shredtools, paper.txt:301-312
We ask whether the 389,234,668 released alignments index on one node, and what a regional query then costs.
% src: data/partners.tsv column 2 summed over 466 files
The 466 TPA files take 25,070,431,026 bytes, and the AGC sequence archive takes 3,295,070,733 bytes \citep{Deorowicz2023agc}.
% src: phase2/s3/tpas.tsv (sum); supporting/plan.md section 4 (local ls of HPRC_r2_assemblies_0.6.1.agc)
By the field sizes in the code, a single-index build holds at least 59,163,669,536 bytes in memory at once.
% src: impg src/main.rs:11621-11760 and src/impg.rs:1657-1735; supporting/plan.md section 4 (389,234,668 x 56 + 778,469,336 x 48)
We therefore build on one cluster node \authornote{node name, CPU model, RAM and scratch bytes, decision D5}.
We sample 1,000 regions each of 1 kb, 10 kb, 100 kb and 1 Mb on non-reference haplotypes, plus 200 held out per size.
% src: supporting/plan.md section 7 E2 (seed 2)
Each region runs as its own impg process on one thread, five times with a warm page cache, at depth 1 and 2.
% src: supporting/plan.md section 7 E2
Each run uses exact and approximate mode under \texttt{/usr/bin/time -v}.
A batch run of the 10 kb set at 1, 4, 16 and 48 threads measures throughput and output identity.
One query covers all of CHM13 chromosome 1.
vg chunk from vg v1.76.1 extracts each region from the GBZ version 3 file of the PGGB graph built from the same alignments \citep{Garrison2018vg,Siren2022gbz,Garrison2024pggb} \authornote{check that vg v1.76.1 reads GBZ version 3, which the release README says needs ``a recent release of vg'' without naming a version}.
% src: HPRC2026resources README line 151 (same all-vs-all alignment), line 159 (v3.gbz requires a recent release of vg); phase2/s3/pggb_wg.tsv (v3.gbz 6,398,170,680 bytes); vg version on the workstation; supporting/plan.md section 14 (v3 GBZ read by vg v1.76.1 not run)
Shredtools extract reads its public HPRC release 2 index, hprcv2\_merged.bumbl, which takes 27,478,774,573 bytes with its \texttt{.bi} file \citep{Shivakumar2026shredtools,Shivakumar2025mumemto}.
% src: phase2/judge-feasibility.md section 1 (aws s3 ls s3://genome-idx/mumemto/hprc/: hprcv2_merged.bumbl 27,424,443,909 + .bumbl.bi 54,330,664 bytes)
All three tools query the same 1,000 GRCh38 regions of 10 kb.
% src: supporting/plan.md section 7 E2
impg and Shredtools return intervals, and vg chunk returns a subgraph, so this comparison measures cost, not output.
Shredtools also answers queries from a web browser through HTTP Range requests on its cloud-hosted index, and impg has no remote access.
% src: paper.txt:25-26, 476, 501, 572 (HTTP Range requests), 640-641 (browser tool); supporting/plan.md section 10 (remote HTTP access not implemented)
impg pays no alignment cost here, because the PGGB build produced the alignments, and the release does not state that cost.
% src: HPRC2026resources README line 151; supporting/plan.md section 12 Q3
The build gave \pending{E2}{outcome, wall and CPU h, peak RSS GB, index bytes} (Fig.~\ref{fig:cost}a).
Per-query cost was \pending{E2}{median and p95 s, peak RSS GB, per size, depth and mode} (Fig.~\ref{fig:cost}b).
On the shared GRCh38 regions, vg chunk took \pending{E2}{median and p95 s, peak RSS GB}.
Shredtools took \pending{E2}{median and p95 s, peak RSS GB} on the same regions.
The whole-chromosome query took \pending{E2}{s and peak RSS GB, or the limit reached}.
Batch output across thread counts was \pending{E2}{identical or not, differing rows} (Fig.~\ref{fig:cost}d).
We make no storage claim.
Exact tracepoint queries need the AGC archive, and GBZ stores the sequences itself.
% src: impg 0987752 src/main.rs:4186-4215 (tracepoint files need --sequence-files unless --approximate with bed or bedpe)
The TPA set and the AGC archive take 28,365,501,759 bytes, 4.433 times the 6,398,170,680 bytes of GBZ version 3, before any impg index (Fig.~\ref{fig:cost}c).
% src: 25,070,431,026 + 3,295,070,733 = 28,365,501,759; 28,365,501,759 / 6,398,170,680 = 4.433 (phase2/s3/tpas.tsv, phase2/s3/pggb_wg.tsv, supporting/plan.md section 4)
As a consistency check, we compare impg intervals at depth 2 with vg chunk path intervals on 1,000 GRCh38 and 1,000 non-reference regions.
% src: supporting/plan.md section 7 E8 (seed 10)
The PGGB graph drops exact matches shorter than 311 bp, so this check measures consistency, not truth \citep{Garrison2023seqwish,Garrison2024pggb}.
% src: HPRC2026resources README line 155 (pggb -p 98 -k 311)
Per-region base-pair Jaccard was \pending{E8}{median and IQR, segmental duplication, satellite and other strata} (Supplementary Fig.~S3).

\begin{figure}[!tpb]
\centering
\fbox{\parbox{0.95\linewidth}{\small Planned panels.
(a) Index build cost of the 466 TPA files.
(b) Per-query wall time and peak memory by region size for impg, vg chunk and Shredtools.
(c) Bytes on disk of released and derived files.
(d) Batch wall time and output identity by thread count.}}
\caption[Cost on the whole release]{\textbf{Cost on the whole release.}
(a) Index build of the 466 HPRC release 2 TPA files in single-index mode.
Bars give wall time (h), CPU time (h), peak resident memory (GB) and index size (bytes), each axis from 0.
Data from E2, \texttt{/usr/bin/time -v} on the node of Supplementary Note S5.
(b) Per-query wall time (s, y axis from 0) against region size (x axis, 1 kb, 10 kb, 100 kb and 1 Mb).
Points are medians and bars are 95th percentiles of 1,000 regions times 5 replicates.
Series are impg exact and approximate at depth 1 and 2, vg chunk on GBZ version 3, and Shredtools extract on hprcv2\_merged.bumbl.
The vg and Shredtools series use the 1,000 GRCh38 regions of 10 kb.
A second row gives peak resident memory (GB, y axis from 0).
Data from E2.
(c) Bytes on disk (y axis from 0) of eight files or sets.
These are the per-target PAF.gz set, the TPA set, the AGC archive (3,295,070,733 bytes), the impg index, GBZ versions 1 and 3, the gbz-base database and the Shredtools index.
Data from the S3 listings of 2026-11-08 and, for the impg index, from E2.
(d) Batch wall time (s, y axis from 0) against threads (x axis, 1, 4, 16 and 48) for 1,000 non-reference 10 kb regions.
Marks flag thread counts whose output differs from the 1-thread output.
Data from E2.}
\label{fig:cost}
\end{figure}

\subsection{Reach and reference dependence}
\label{sec:res-reach}

Every assembly in the release aligns to GRCh38 and CHM13, whose two files list all 465 other haplotypes \citep{Schneider2017grch38,Nurk2022chm13}.
% src: data/partner_summary.txt (refs 465 each); data/pairs_out.txt
Between non-reference haplotypes, the index walks 22,809 of 107,416 unordered pairs (21.2343\%).
% src: data/pairs_out.txt
At the level of pairs, every non-reference haplotype reaches all 463 others within two hops without either reference (Fig.~\ref{fig:release}c).
% src: data/pairs_out.txt (2 hops min 463 median 463); sections/04.release.tex fig:release panel c
Pair-level reach does not measure bases, so we ask how many bases a query reaches without the two references.
We also ask how well transitive projection recovers a held-out direct alignment.
In the first test we project 1,000 random 1 kb GRCh38 regions to CHM13, plus 200 held out.
% src: supporting/plan.md section 7 E3a (seed 3)
One index holds the GRCh38 and CHM13 target files, and a second index holds them without the direct GRCh38 and CHM13 records.
% src: supporting/plan.md section 7 E3a
The second index reaches CHM13 only through other haplotypes, at depth 2 and 3.
We score it against the direct projection and against a public GRCh38-to-CHM13 chain file \citep{Kent2003chain} \authornote{choose the chain file and record its provenance}.
In the second test we choose 10 non-reference haplotypes, each with one partner whose alignments only the haplotype's own file stores.
% src: supporting/plan.md section 7 E3b (seed 4); data/pairs_out.txt (2,966 of 25,221 ordered pairs are stored in both files)
We replace each such file by its per-target PAF without the partner's records, so no direct alignment links the two.
% src: supporting/plan.md section 7 E3b
We query 100 random 10 kb regions per haplotype, 20 of them held out, at depth 1, 2, 3, 4 and unbounded.
% src: supporting/plan.md section 7 E3b
Each depth runs with the references and without them, through a subset list that removes GRCh38 and CHM13 as outputs and as bridges.
% src: impg src/impg.rs:2496-2505 (subset filter applied before a projection joins the next frontier); VERIFY_CORRECTIONS.md:18
The truth is the direct projection of the same regions on the full index.
At depth 1 no direct alignment remains, so this depth serves as a control.
Metrics are bp recall, bp precision, the exact-boundary rate at 0, 1, 10 and 100 bp, and bases reached per depth.
% src: supporting/plan.md section 7 E3
Strand agreement waits for the strand decision, because impg sets BED strand per hop and does not compose it back to the query.
% src: impg src/impg.rs:282-286 and src/main.rs:12072-12079; supporting/plan.md sections 9 C7 and 13 D3
GRCh38 to CHM13 projection through haplotypes gave \pending{E3}{bp recall and precision against direct projection and against the chain, depth 2 and 3}.
Held-out partner recall was \pending{E3}{median per-region bp recall with and without references, per depth} (Fig.~\ref{fig:recall}c).
Bases reached were \pending{E3}{bp per depth, with and without references} \authornote{Fig.~\ref{fig:release} has panels a to c only. Choose the panel that shows bases reached per depth}.
We report unbounded runs that do not finish with their limit.

\subsection{Sparse alignment sets and transitive error}
\label{sec:res-sparse}

The number of genome pairs grows as $n(n-1)/2$, so 90 haplotypes give 4,005 pairs and 142 strains give 10,011.
We ask how much homology transitive projection recovers, by depth, as the fraction of aligned pairs falls.
The tuning set is the all-pairs chr19 alignment of 90 HPRC year-1 haplotypes from Section~\ref{sec:res-genes} \citep{Liao2023hprc}.
% src: supporting/plan.md section 7 E4 and E5
The holdout is the ScRAP set of 142 \textit{Saccharomyces cerevisiae} assemblies, which never sets a parameter \citep{ODonnell2023scrap}.
% src: supporting/plan.md section 7 E4; ODonnell2023scrap
We align all pairs with wfmash through \texttt{impg align} and keep its default filters \citep{Guarracino2024wfmash,Garrison2026sweepga}.
% src: impg align --help (--sparsify none default, plane-sweep and scaffold filters on by default)
From each all-pairs set we draw four subsets of pairs.
Two are random 6\% and 17\% subsets, the two pair fractions of the Shredtools benchmark \citep{Shivakumar2026shredtools}.
% src: paper.txt:301-303 (6% of pairwise alignments), 310-311 (17%, the default auto setting); the preprint gives no selection rule for either set
The third holds the pairs that \texttt{-{}-sparsify auto} selects, and the fourth is a star to one genome fixed in advance.
% src: supporting/plan.md section 7 E4; section 13 D8 (hub strain)
Above 50 haplotypes, \texttt{auto} keeps the 5 nearest and 2 farthest Mash neighbours of each haplotype plus a 0.05 random fraction of pairs \citep{Ondov2016mash}.
% src: sweepga 4713ce0 src/knn_graph.rs:514-532 (auto rule) and 243-280 (tree pairs)
impg groups FASTA records by PanSN haplotype before this step, and it always aligns contigs within one haplotype.
% src: impg src/commands/align.rs:147-179; sweepga 4713ce0 src/knn_graph.rs:589-618 (within-haplotype pairs always added), 636-677 (haplotype-level selection)
We draw the random subsets in the harness with a fixed seed, because \texttt{random:<f>} keeps each pair by a hash threshold.
% src: sweepga 4713ce0 src/knn_graph.rs:371-392
We record whether the \texttt{auto} pair count at commit 0987752 matches the 17\% that Shredtools called the default auto setting of impg v0.4.0.
% src: paper.txt:310-311 (default auto setting), 838-839 (impg v0.4.0); sweepga 4713ce0 src/knn_graph.rs:514-532, pinned in Cargo.lock at 0987752
The truth is direct projection on the all-pairs index.
We project 1,000 random 10 kb regions per set at depth 1, 2, 3, 4 and unbounded on each subset index.
% src: supporting/plan.md section 7 E4 (seed 5)
Metrics are bp recall, bp precision, the exact-boundary rate at 0, 10 and 100 bp, and haplotypes reached.
We also report pairs aligned and alignment CPU time per subset.
On chr19, recall at depth 2 was \pending{E4}{median per-region bp recall per subset} (Fig.~\ref{fig:recall}a).
On ScRAP, recall at depth 2 was \pending{E4}{median per-region bp recall per subset}.
Exact boundaries within 10 bp occurred at \pending{E4}{rate per subset and depth, chr19 and ScRAP}.
Alignment cost was \pending{E4}{pairs aligned and CPU hours per subset, against all pairs}.
The \texttt{auto} rule selected \pending{E4}{haplotype pairs and percentage at n = 90 and n = 142}.

\subsection{Gene-level accuracy and Shredtools on the same input}
\label{sec:res-genes}

We ask how often each tool returns the annotated copy of a gene in each haplotype, and at what cost.
Shredtools compared with impg on 92 HPRC year-1 chr19 assemblies, and the preprint does not list them \citep{Shivakumar2026shredtools}.
% src: paper.txt:283-312
We use the paths of the 90 haplotypes in the HPRC year-1 PGGB chr19 graph, built with \texttt{-n 90} \citep{Liao2023hprc,Garrison2024pggb}.
% src: phase2/judge-feasibility.md section 1 (freeze1 README); S3 listing of freeze1/pggb/chroms (chr19.hprc-v1.0-pggb.gfa.gz, 345,202,999 bytes)
The truth is the CAT annotation of GENCODE v38 on HPRC year-1 assemblies, 89 files from its run on CHM13 \citep{Fiddes2018cat,Frankish2021gencode}.
% src: aws s3 ls of submissions/FC7E9302-5A35-4A4F-9D4A-B3E957118E36--Y1_CAT_ANNOTATION_GENCODE38 on 2026-11-09 (178 files, 89 named *.chm13.gff3.gz)
Queries are the CHM13 chr19 gene intervals of the same run, scored on haplotypes with both a path and a CAT file.
CAT projects genes through whole-genome alignments, so this truth is not independent of alignment \citep{Fiddes2018cat}.
We align all 4,005 haplotype pairs with wfmash through \texttt{impg align}.
From these we draw random 17\% and 6\% subsets, the \texttt{auto} pairs and a CHM13 star.
% src: supporting/plan.md section 7 E5
Shredtools runs with Mumemto v1.4.0 on the same per-haplotype FASTA files \citep{Shivakumar2025mumemto} \authornote{Shredtools commit to pin}.
% src: paper.txt:637-640 (Mumemto v1.4.0; code at github.com/vikshiv/shredtools)
We time alignment at 48 threads and single queries at 1 thread.
We time the index build at 1, 8, 16 and 48 threads.
We record the peak memory of each stage.
% src: sections/08b.supplement-experiments.tex Table S1 E5 block (align -t 48, index t in 1 8 16 48, query -t 1); supporting/plan.md section 7 E5; paper.txt:346-349 (48 threads)
Metrics per gene and haplotype are containment of the CAT interval, flank excess in bp, wrong-locus rate and copies returned.
We also time 1,000 random 10 kb regions per tool.
Shredtools reported that its index build was 86 times faster than impg at 48 threads, a rounded factor, and the impg time included wfmash alignment \citep{Shivakumar2026shredtools}.
% src: paper.txt:299-300 (printed with a tilde); paper.txt:345-348 (Fig. 2 legend: impg index using wfmash)
It reported 0.484 s on average and 79 MB of peak memory per query, against 4.01 s and 1.66 GB for impg on all pairs.
% src: paper.txt:308-311; impg version v0.4.0 from its reference list, paper.txt:838-839
On 17\% of pairs it reported 0.21 s and 12 MB for impg.
% src: paper.txt:311-312
These are their numbers for impg v0.4.0, not our results.
% src: paper.txt:838-839
Containment was \pending{E5}{median per gene and haplotype, impg per subset and Shredtools} (Fig.~\ref{fig:recall}b).
The wrong-locus rate was \pending{E5}{fraction of gene and haplotype pairs, per tool}.
Flank excess was \pending{E5}{median bp per tool}.
Stage costs were \pending{E5}{wall time in h and peak RSS in GB for alignment, index and query, per tool and thread count}.
Per-query cost was \pending{E5}{median and p95 wall time in s, peak RSS in MB, per tool}.

\begin{figure}[!tpb]
\centering
\fbox{\parbox{0.95\linewidth}{\small Planned panels.
(a) Recall and exact boundaries against the fraction of aligned pairs.
(b) Gene-copy accuracy of impg subsets and Shredtools against CAT.
(c) Held-out partner recall on HPRC release 2 with and without the reference bridges.}}
\caption[What transitive projection recovers]{\textbf{What transitive projection recovers.}
(a) Base-pair recall (y axis, 0 to 1) against the fraction of aligned genome pairs (x axis, 0 to 1).
Subsets are random 6\%, random 17\%, the \texttt{-{}-sparsify auto} pairs and a star, with one line per depth (1, 2, 3, 4 and unbounded).
A second row gives the exact-boundary rate within 10 bp (y axis, 0 to 1).
The left column shows the chr19 tuning set of 90 HPRC year-1 haplotypes, and the right column the ScRAP 142 holdout.
Data from E4.
(b) Per gene and haplotype containment of the CAT GENCODE v38 interval (y axis, 0 to 1), flank excess (bp) and wrong-locus rate.
Series are impg on all pairs, 17\%, 6\%, the \texttt{auto} pairs and a CHM13 star, and Shredtools.
Data from E5.
(c) Held-out partner bp recall (y axis, 0 to 1) against depth (x axis, 1 to 4 and unbounded), with and without reference bridges.
Points are 10 query haplotypes of HPRC release 2 (E3b).
An inset gives GRCh38 to CHM13 recall through haplotypes against direct projection and the public chain (E3a).}
\label{fig:recall}
\end{figure}

\subsection{Local graphs and the restriction of the implied graph}
\label{sec:res-local}

\texttt{impg query -o gfa} builds a graph from the homologous intervals that a query finds.
It keeps the merged intervals and re-aligns them with the chosen engine, so its graph is not cut from a global graph.
% src: impg src/main.rs:12326-12340; src/lib.rs:926-951
Proposition~\ref{prop:restrict} gives a condition for an exact cut.
Seqwish on alignments clipped to a closed set gives the global seqwish node classes on that set, if seqwish \texttt{-k} is at most 1 \citep{Garrison2023seqwish}.
% src: supporting/opportunities.md section 1.6 (Proposition 2)
We test this on real alignments and measure how far the default closure and the \texttt{-o gfa} path lie from it.
The tuning set is the yeast chrV fixture of impg, 7 sequences and 4,029,392 bp, with 100 random 10 kb loci.
% src: impg tests/test_data/yeast.chrV.fa.gz (7 records, 4,029,392 bp; VERIFY_CORRECTIONS.md:100); supporting/plan.md section 7 E6 (seed 7)
The holdout is the C4 region of GRCh38 chr6 in 20 HPRC release 2 haplotypes.
% src: supporting/plan.md section 7 E6 (GRCh38#0#chr6:31891045-32123783, seed 8)
We align all pairs with wfmash and induce the global graph with seqwish v0.7.11-35-gaab5bee at \texttt{-k} 0 and 23.
% src: seqwish --version and --help on the workstation (-k default 0)
For each locus, an unbounded closure with minimum hop length 0, snapping 0 and no merging writes the clipped alignments as PAF rows with CIGARs.
% src: impg src/main.rs:12240-12295 (PAF rows carry cg:Z); supporting/opportunities.md section 1.6
Seqwish on these rows gives the restricted graph, and we compare node classes base by base on the closure.
Two more arms use the closure at the command-line defaults and the \texttt{-o gfa} seqwish engine.
% src: impg query --help (-m 2, --min-transitive-len 101, --min-distance-between-ranges 10)
The defaults are depth 2, a minimum hop of 101 bp and 10 bp snapping.
% src: impg query --help; impg src/main.rs:4253-4255
The metric is the number and fraction of closure bases whose node class differs from the global graph.
At \texttt{-k} 0 the restricted graph differed on \pending{E6}{bases, count and fraction, yeast loci and C4} (Fig.~\ref{fig:local}a).
At \texttt{-k} 23, at the defaults and through \texttt{-o gfa}, the counts were \pending{E6}{bases and fraction per arm}.
Graphs from \texttt{-o gfa} remain graphs built from the homologous intervals, because that path re-aligns them.
% src: impg src/main.rs:12326-12340

\subsection{Locus refinement}
\label{sec:res-refine}

\texttt{impg refine} extends a locus to maximize the number of sequences, samples or haplotypes that span both of its ends, and counts sequences by default.
% src: impg refine --help; impg 0987752 src/main.rs:7932-7934 (pansn_mode defaults to PanSnLevel::Sequence); VERIFY_CORRECTIONS.md:64
COSIGT uses \texttt{impg query} and \texttt{impg refine} to extract loci before genotyping \citep{Bolognini2026cosigt}.
% src: Bolognini2026cosigt; supporting/plan.md section 5.1
Refine has one unit test and no command-line test.
% src: VERIFY_CORRECTIONS.md:91 (refine 1 unit test); supporting/plan.md section 9 C11
We run it at TERT and IGKV6-21, the loci of two Shredtools examples \citep{Shivakumar2026shredtools}.
% src: paper.txt:443-470
At IGKV6-21, Shredtools first found bounds of 328,599 bp and 74,147 bp across 15 T2T genomes.
% src: paper.txt:445-448
It then dropped two genomes from cell lines with VDJ deletions and extracted the gene with 0 bp bounds from the other 13.
% src: paper.txt:448-454
At TERT it reported bounds of 25 to 35 kbp on either side before enhance and exact bounds across all 476 genomes after it.
% src: paper.txt:463-470
Refine changes the interval, while zoom changes the set of genomes.
We run refine on the E2 index with haplotype counting, 1,000 bp end spans, 1,000 bp steps and extension up to half the locus.
% src: impg refine --help (--pansn-mode, --span-bp 1000, --extension-step 1000, --max-extension 0.5)
We run it on all 466 files and on the T2T haplotypes among them, and Shredtools runs enhance and zoom on its own index.
% src: supporting/plan.md section 7 E10
Refine reached \pending{E10}{haplotypes spanning both ends before and after, extension bp per side, wall time in s, per locus} (Fig.~\ref{fig:local}b).
Zoom kept \pending{E10}{genomes kept and bounds in bp, per locus}.

\begin{figure}[!tpb]
\centering
\fbox{\parbox{0.95\linewidth}{\small Planned panels.
(a) Bases whose node class differs between restricted and global seqwish graphs, per arm.
(b) Haplotypes spanning both ends against interval length at TERT and IGKV6-21.}}
\caption[Local graphs and locus boundaries]{\textbf{Local graphs and locus boundaries.}
(a) Fraction of closure bases whose node class differs between the restricted and the global seqwish graph (y axis from 0).
The x axis has four arms: seqwish \texttt{-k} 0, seqwish \texttt{-k} 23, the closure at command-line defaults, and the \texttt{-o gfa} seqwish engine.
Points are the 100 yeast chrV loci and the C4 holdout in 20 haplotypes.
Data from E6.
(b) Haplotypes spanning both ends (y axis from 0) against interval length (x axis, bp) at TERT and IGKV6-21.
Series are impg before and after refine, and Shredtools extract before and after enhance and zoom.
Data from E10.}
\label{fig:local}
\end{figure}

\subsection{Partition work units on yeast}
\label{sec:res-partition}

\texttt{impg partition} grows partitions from windows of one sequence and masks each assigned base, so later partitions skip it.
% src: VERIFY_CORRECTIONS.md:93 (cover by masking); impg src/commands/partition.rs:250-256
We ask whether this covers every base of an alignment-backed pangenome exactly once.
We partition the ScRAP all-pairs index with 10 kb windows, 1 kb merging and haplotype selection, fixed before the run.
% src: supporting/plan.md section 7 E9 (-w 10000 -d 1000 --selection-mode haplotype)
An audit counts partitions in the BED files, because the logged count can exceed them after rehoming.
% src: impg src/commands/partition.rs:142-156, 676-678, 746-748; VERIFY_CORRECTIONS.md:88
It also reports uncovered bases, bases owned more than once and partitions that join different S288C chromosomes.
The partition gave \pending{E9}{partitions, member intervals, uncovered bp, bp owned twice, cross-chromosome partitions, wall time, peak RSS} (Supplementary Fig.~S4).
